\chapter{Robustness to Perturbation and Prompt Rewording: When a Stable Answer Hides a Drifted Explanation (Day 9)}
\label{chap:day09}

\section{Motivation}

Robustness is one of the six metrics defined for evaluating black-box explainable-AI frameworks
\citep{arreche2024exai}: does a judgment, and the explanation behind it, survive small, realistic
perturbations to the input, or does it depend brittlely on exact pixel values and exact wording?
Chapters 3--8 each established a different facet of how Florence-2 behaves on a single, fixed input
--- its baseline answers, its grounded regions, its masking sensitivity, its attention concentration,
its stochastic-decoding stability, its computational cost --- but none of them perturbed the
\emph{input itself} and asked whether the answer and the explanation move together, move separately,
or do not move at all. Day 9's task was to answer exactly that question, along two axes: realistic
image perturbations (blur, lighting, occlusion, contrast) and prompt rewording (asking the same
safety question in different words). The result is the richest chapter in this pilot, both in volume
of findings and in the weight one of them carries: \textbf{a stable final answer can sit directly on
top of an explanation that has completely drifted}, a result that reframes how every other chapter's
numbers should be read.

\section{Implementation}

\texttt{perturbations.py} implements four Level-1 image perturbations and one stretch-test
compositing operation. \textbf{\texttt{blur}} applies Gaussian blur with \texttt{ksize=8.0}.
\textbf{\texttt{low\_light}} applies gamma correction with \texttt{gamma=2.5}, darkening the image.
\textbf{\texttt{occlude}} draws a centered black square covering \texttt{frac=0.2} of the image's
area. \textbf{\texttt{contrast\_shift}} scales pixel values around the mid-point by
\texttt{factor=0.3}. \textbf{\texttt{paste\_patch}} composites an arbitrary patch image onto an
arbitrary box and is used only by the stretch test described in
Section~\ref{sec:day09-stretch}, not by the four Level-1 perturbations.

\texttt{robustness.py}'s \texttt{RobustnessResult} dataclass records four fields per rerun:
\texttt{answer\_changed} (did the final yes/no safety-rule answer flip), \texttt{confidence\_drop}
(baseline confidence minus perturbed confidence; a negative value means confidence \emph{rose}),
\texttt{object\_box\_iou} (intersection-over-union between the baseline's and the rerun's safety-object
box, \texttt{NaN} if either side has no object box to compare), and \texttt{worker\_lost} --- a flag
that fires specifically when the baseline run detected at least one worker but the perturbed rerun
detected none. \texttt{worker\_lost} exists because Chapter~\ref{chap:day05} already found that
masking the only detected worker reroutes \texttt{\_answer\_presence\_rule} into a degenerate
scene-level fallback branch, flipping the answer for reasons disconnected from the safety concept
actually being tested. The same risk applies to any perturbation that happens to make the worker
undetectable, not only to a deliberate mask, so this chapter inherited Chapter~5's diagnostic rather
than reinventing it.

\texttt{scripts/09\_robustness\_test.py} ran three test families over the 163-sample pilot. \textbf{Level
1} reruns all 163 samples through each of the four image perturbations above, 652 reruns in total.
\textbf{Level 2} reruns 138 of the 163 samples with the rule's second phrasing instead of an image
change (\texttt{phrasing\_index=1}); rule\_4 is excluded, because \texttt{answer\_rule} routes rule\_4
to \texttt{\_answer\_rule\_4}, which has no \texttt{phrasing\_index} parameter at all --- rule\_4 is a
fixed proximity pair (worker, excavator), not a phrase list, so passing \texttt{phrasing\_index=1} for
it would have silently been a no-op rerun of the exact same query, not a real test of wording
sensitivity. This was caught before any aggregation code was written and rule\_4's 25 excluded rows
are logged with \texttt{excluded\_reason="rule\_4\_has\_no\_second\_phrasing"} rather than silently
scored as 0\% change, which would have understated rule\_4's true (untested) rewording sensitivity by
fabricating a result for it. \textbf{The stretch test} is a 20-sample synthetic hard-hat-patch
experiment, covered separately in Section~\ref{sec:day09-stretch}. All Level-1 and Level-2 reruns use
the \textbf{same deterministic decoding as the logged baseline} (\texttt{num\_beams=3}, no sampling),
a deliberate methodological choice: without it, any answer change observed here could be attributed
either to the perturbation or to Chapter~\ref{chap:day07}'s already-documented sampling noise, and
there would be no way to tell which. Holding decoding fixed isolates perturbation sensitivity as the
only thing that can change between baseline and rerun.

\section{Architecture: A Perturbation Layer}

Figure~\ref{fig:diagram-e} lays out the three parallel branches \texttt{09\_robustness\_test.py}
runs over each baseline sample: a Level-1 image-perturbation branch, a Level-2 reworded-prompt
branch that leaves the image untouched, and a patch-stretch branch that composites a synthetic
hard-hat patch onto the image before rerunning. All three branches converge on the same comparison
step against the logged baseline, computing \texttt{answer\_changed}, \texttt{confidence\_drop},
\texttt{object\_box\_iou}, and \texttt{worker\_lost} uniformly regardless of which branch produced
the rerun.

\begin{figure}[h]
\centering
\resizebox{\textwidth}{!}{%
\begin{tikzpicture}[node distance=1.6cm,
    block/.append style={minimum width=3.6cm, minimum height=1.3cm}]
\node[block] (img) {Original image};
\node[block, below=2.0cm of img, xshift=-5.2cm] (l1) {blur / low\_light /\\occlude / contrast\_shift};
\node[block, below=2.0cm of img] (l2) {phrasing\_index=1\\(no image change)};
\node[block, below=2.0cm of img, xshift=5.2cm] (l3) {paste\_patch\\(head\_box, hard hat)};
\node[block, below=1.8cm of l1] (rerun1) {answer\_rule rerun};
\node[block, below=1.8cm of l2] (rerun2) {answer\_rule rerun};
\node[block, below=1.8cm of l3] (rerun3) {answer\_rule rerun};
\node[block, below=1.8cm of rerun2, minimum width=8.6cm, fill=rulegray] (cmp) {Compare to baseline:\\answer\_changed, confidence\_drop,\\object\_box\_iou, worker\_lost};
\draw[arrow] (img) -- (l1);
\draw[arrow] (img) -- (l2);
\draw[arrow] (img) -- (l3);
\draw[arrow] (l1) -- (rerun1);
\draw[arrow] (l2) -- (rerun2);
\draw[arrow] (l3) -- (rerun3);
\draw[arrow] (rerun1) -- (cmp);
\draw[arrow] (rerun2) -- (cmp);
\draw[arrow] (rerun3) -- (cmp);
\end{tikzpicture}%
}
\caption{Perturbation layer: three parallel robustness branches, Level 1/Level 2/patch-stretch.}
\label{fig:diagram-e}
\end{figure}

Note that Level 2's branch is the only one of the three that leaves the image itself unperturbed ---
the input change is entirely in the text prompt, not the pixels --- which is why it sits in the
middle of the diagram with no image transformation drawn inside its top-row node, only the
\texttt{phrasing\_index=1} relabeling. The patch-stretch branch (right) is structurally a
Level-1-style image perturbation but is kept visually and analytically separate because, unlike
blur/low\_light/occlude/contrast\_shift, it targets one specific rule (rule\_1) and one specific
region (the approximate head box) by construction, rather than perturbing the whole frame uniformly.

\section{Results}

Table~\ref{tab:day09-headline} reports the headline \texttt{answer\_changed} rate for each of the
five perturbation types, alongside the three diagnostic columns that distinguish \emph{why} an
answer changed from merely \emph{whether} it changed.

\begin{center}
\begin{tabular}{p{3.0cm}p{2.0cm}p{1.7cm}p{3.2cm}p{2.1cm}}
\toprule
\textbf{Perturbation} & \textbf{answer\_\allowbreak changed} & \textbf{worker\_\allowbreak lost} & \textbf{mean confidence\_\allowbreak drop} & \textbf{mean object\_\allowbreak box\_iou} \\
\midrule
occlude & \textbf{28.2\%} & 8.6\% & $-0.012$ (rose) & 0.530 \\
reworded\_\allowbreak prompt & 27.5\% & 0.0\% & $-0.053$ (rose) & 0.333 \\
blur & 27.0\% & 0.0\% & 0.013 & 0.287 \\
low\_light & 11.7\% & 0.6\% & 0.011 & 0.816 \\
contrast\_\allowbreak shift & 10.4\% & 0.0\% & 0.008 & 0.825 \\
\bottomrule
\end{tabular}
\captionof{table}{\texttt{answer\_changed} rate by perturbation, with the three diagnostic columns
that explain \emph{why}: \texttt{worker\_lost} (the Chapter-5 fallback-rerouting risk),
mean \texttt{confidence\_drop} (negative means confidence rose), and mean \texttt{object\_box\_iou}
(lower means the object box moved more).}
\label{tab:day09-headline}
\end{center}

Brightness and contrast changes are the mildest perturbations: object boxes barely move
(\texttt{object\_box\_iou} above 0.8 for both) and answers rarely flip (10--12\%). Blur, occlusion,
and prompt rewording are all substantially more disruptive --- answer-change rates between 27\% and
28\%, roughly $2.5\times$ low\_light and contrast\_shift's rates --- but, as Findings 1--4 below show
in turn, \textbf{for three different reasons, not one}. Treating "blur, occlude, and reworded\_prompt
are all about equally disruptive" as a single finding would flatten three structurally different
phenomena into one number; the rest of this chapter is about pulling them back apart.

\section{Finding 1: Occlusion's Fallback Artifact}
\label{sec:day09-finding1}

Occlude's 28.2\% answer-change rate is not uniformly "genuine sensitivity to a covered region."
\textbf{Of occlude's 46 answer-flip cases, 9 (20\%) also have \texttt{worker\_lost=True}} --- the
occluded rerun's fresh detection found no worker at all, routing \texttt{\_answer\_presence\_rule}
into its scene-level fallback branch for reasons disconnected from whether the safety object itself
was actually still visible. The \texttt{worker\_lost} rate under occlude is also far from uniform
across rules: rule\_2 19.2\%, rule\_4 16.0\%, rule\_3 8.2\%, rule\_1 1.6\% --- a more than
tenfold spread between the highest and lowest rule, all under the identical occlusion transform.

\textbf{Traced concretely on image \texttt{0000341} (rule\_3, baseline ``compliant'').}
Figure~\ref{fig:day09-finding1} shows the before/after pair. Before occlusion, \textbf{both baseline
detections were already wrong}: the box labeled "worker" is actually a false positive on the
excavator's cab window, and the box answering the guardrail query is the entire excavator
silhouette, not a guardrail at all. Both boxes happened to satisfy the geometric coverage check the
rule applies, producing a "compliant" answer that was correct only by accident. After occlusion,
\textbf{both} detections vanish --- \texttt{worker\_lost=True}, \texttt{object\_box\_iou=NaN} ---
even though the occluded square does not fully cover either of the original (wrong) boxes.
Occlusion is disrupting Florence-2's whole-scene grounding context, not merely blacking out the
literal pixels the two boxes occupied. With no object box surviving the rerun, the fallback branch
defaults to "violation" purely because no object box was found, flipping
\texttt{answer\_changed=True}.

\begin{figure}[h]
\centering
\begin{subcaptionblock}{0.48\linewidth}
\centering
\includegraphics[width=\linewidth]{robustness/verify/0000341_before.png}
\caption{Before: both the "worker" box (a false positive on the excavator's cab window) and the
guardrail-query box (the entire excavator silhouette) are wrong, but geometrically satisfy the
coverage check.}
\end{subcaptionblock}
\hfill
\begin{subcaptionblock}{0.48\linewidth}
\centering
\includegraphics[width=\linewidth]{robustness/verify/0000341_after_occlude.png}
\caption{After occlude: both detections vanish (\texttt{worker\_lost=True},
\texttt{object\_box\_iou=NaN}), even though the occluded square does not fully cover either
original box.}
\end{subcaptionblock}
\caption{Image \texttt{0000341}, rule~3: occlusion's fallback artifact, traced. The flip to
"violation" reflects a degenerate fallback branch firing because no object box survived, not a
correct detection that the guardrail became occluded.}
\label{fig:day09-finding1}
\end{figure}

Read in an aggregate table alone, this case is indistinguishable from "the model correctly noticed
the safety equipment became occluded." \textbf{It is not.} It is the same degenerate-fallback
mechanism Chapter~\ref{chap:day05} already flagged for deliberate masking, here triggered by a
perturbation instead. The two settings differ only in what disabled the detection --- a hand-drawn
mask there, a black occlusion square here --- not in the underlying failure mode, which is that
\texttt{\_answer\_presence\_rule}'s fallback branch answers from the mere absence of any object box,
with no way to distinguish "the safety object is genuinely gone" from "the detector lost its grip on
the whole scene." Any future Bounded Completeness analysis (Chapter~\ref{chap:day10}) needs to apply
this same \texttt{worker\_lost} check rather than reporting a raw \texttt{no\_usable\_explanation}
rate blindly, since some fraction of "no usable explanation" cases will be this same fallback
artifact rather than a genuine absence of grounding.

\section{Finding 2: A Stable Answer Hiding a Drifted Explanation}
\label{sec:day09-finding2}

\textbf{This is the central result of Chapter~9, and arguably of the entire pilot.} Quantified
across all 611 stable-answer rows that have a comparable object box on both sides (baseline and
rerun both detected at least one object box, and the answer itself did not change):
\textbf{28.3\% have \texttt{object\_box\_iou < 0.3}} --- the answer stayed exactly the same, but the
box the model points to as evidence moved to something close to unrelated. This is wildly uneven by
perturbation, as Table~\ref{tab:day09-finding2} shows.

\begin{center}
\begin{tabular}{lr}
\toprule
\textbf{Perturbation} & \textbf{\% of stable answers with \texttt{object\_box\_iou < 0.3}} \\
\midrule
blur & \textbf{58.8\%} \\
reworded\_prompt & 56.7\% \\
occlude & 32.4\% \\
low\_light & 5.7\% \\
contrast\_shift & 2.8\% \\
\bottomrule
\end{tabular}
\captionof{table}{Of samples whose final answer did not change, the percentage whose evidence box
nonetheless drifted to \texttt{object\_box\_iou < 0.3}, by perturbation.}
\label{tab:day09-finding2}
\end{center}

Under blur and reworded\_prompt, \textbf{more than half} of all stable-answer samples are hiding a
near-total relocation of the evidence box. This is not a tail effect or a rare edge case; it is the
majority behavior for two of the five perturbation types.

\textbf{Traced concretely on image \texttt{0000034} (rule\_2, baseline ``compliant'',
\texttt{answer\_changed=False}).} Figure~\ref{fig:day09-finding2} shows the before/after pair: the
single best image in this pilot for demonstrating that answer-level robustness alone is not enough.
The worker's entire baseline box sits \textbf{inside} the occluded square. A reasonable expectation
would be that occluding the worker entirely either loses the detection (\texttt{worker\_lost=True})
or, at minimum, forces the evidence box somewhere conspicuously different in a way that shows up as
an obvious failure. Neither happens cleanly: the rerun still reports \texttt{worker\_lost=False},
because it found \emph{a} worker box somewhere in the image, and the final answer is unchanged ---
yet \texttt{object\_box\_iou=0.005}. The object box moved to a near-completely different location in
the frame, while every answer-facing signal --- the final yes/no answer, the \texttt{worker\_lost}
flag --- reports business as usual.

\begin{figure}[h]
\centering
\begin{subcaptionblock}{0.48\linewidth}
\centering
\includegraphics[width=\linewidth]{robustness/verify/0000034_before.png}
\caption{Before: the worker's baseline box, fully visible.}
\end{subcaptionblock}
\hfill
\begin{subcaptionblock}{0.48\linewidth}
\centering
\includegraphics[width=\linewidth]{robustness/verify/0000034_after_occlude.png}
\caption{After occlude: the worker's entire baseline box sits inside the occluded square; the
rerun still finds \emph{a} worker box (\texttt{worker\_lost=False}) and the answer is unchanged ---
but \texttt{object\_box\_iou=0.005}, an almost total relocation.}
\end{subcaptionblock}
\caption{Image \texttt{0000034}, rule~2: \textbf{the centerpiece finding of this chapter.} An
answer-level robustness number alone would call this sample perfectly robust. It is not, at the
explanation level.}
\label{fig:day09-finding2}
\end{figure}

\textbf{An answer-level robustness number alone would call this sample perfectly robust; it is not,
at the explanation level.} This is the strongest single result in the pilot motivating the report of
answer-level and explanation-level robustness as \textbf{two separate numbers, never blended into
one}. A single blended robustness score computed only from \texttt{answer\_changed} would report
\texttt{0000034} as a clean, unremarkable success --- and would do so for 28.3\% of all stable-answer
rows across the dataset, systematically certifying as "robust" a population of samples whose
explanations have, in fact, drifted to evidence the model is no longer actually looking at. This
mirrors Chapter~\ref{chap:day06}'s decision to report visual and text attribution as two separate,
never-merged streams: here, the two streams that must not be merged are what the model
\emph{answers} and what the model \emph{points to as why}, and Finding~2 is the directly-traced,
concrete failure case that motivates keeping them apart, rather than an assertion of that separation
on principle alone.

\section{Finding 3: Not All Phrasings Are Synonyms}

Level-2 answer-change rate, broken down by rule, surfaces a pattern that initially looks like
rewording sensitivity but is better read as a prompt-design issue specific to one rule:

\begin{center}
\begin{tabular}{lr}
\toprule
\textbf{Rule} & \textbf{reworded\_prompt answer\_changed} \\
\midrule
rule\_2 (harness / lanyard) & \textbf{50.0\%} \\
rule\_1 (hard hat / vest) & 31.7\% \\
rule\_3 (guardrail / edge barrier) & 10.2\% \\
\bottomrule
\end{tabular}
\end{center}

rule\_3's two phrasings ("guardrail" versus "edge protection barrier") are near-synonymous, and
Florence-2 grounds them consistently --- only a 10.2\% answer-change rate. rule\_2's two phrasings
are \textbf{not} really synonyms in the way rule\_3's are: a safety harness is a body garment worn
by the worker, and a fall-protection lanyard is the tether attached to that garment. These are two
different physical objects that happen to co-occur in the same PPE concept, not one object referred
to by two interchangeable names. Asking Florence-2 to ground "harness" versus "lanyard" can
legitimately produce two different bounding boxes on the same scene without either grounding being
wrong --- the model is correctly answering two slightly different questions, not failing to
recognize that one question was reworded. This is a prompt-design caveat for future revisions of
\texttt{prompts.py}, not a model robustness failure: rule\_2's two "phrasings" should probably be
treated as two distinct concepts rather than as wording variants of a single concept, the next time
\texttt{RULE\_QUERIES} is revised.

\section{Finding 4: An Uncalibrated Confidence Proxy}

Mean \texttt{confidence\_drop} is \emph{negative} for both occlude ($-0.012$) and reworded\_prompt
($-0.053$) --- average confidence went \textbf{up}, not down, after the perturbation, despite both
perturbations having among the highest answer-change rates in Table~\ref{tab:day09-headline}. If
confidence behaved as a correctness signal, the perturbations that most disrupt the answer should be
the ones that most erode confidence; instead, the two most disruptive perturbations are associated
with the model growing \emph{more} confident on average. This confirms, again, what
Chapters~\ref{chap:day03} and~\ref{chap:day05} already cautioned: Florence-2's reported "confidence"
is a mean token-generation-probability proxy, not a calibrated correctness signal. Removing visual
information (occlusion) or changing the query wording can make the model surer of a different, not
necessarily more correct, detection --- confidence tracks how decisively the decoder committed to
\emph{some} answer, not whether that answer is right.

\section{Stretch Test: A Synthetic Hard-Hat Patch}
\label{sec:day09-stretch}

The stretch test pastes a synthetic hard-hat patch over the approximate head region of 20
\texttt{ppe\_violation}/rule\_1 candidates and asks whether the patch flips the answer from
"violation" to "compliant" --- a small, bounded probe of whether Florence-2's grounding can be
spoofed by a crude, synthetic visual cue, scoped deliberately small per
\texttt{technical-limitations-take-and-mitigation.md}'s guidance that a "Level 3" spoofing test
should be a small part of this pilot, not a full promise of a demonstrated vulnerability.

\textbf{The denominator-correction story.} The naive headline number is a 15\% (3/20) flip rate.
That number was caught and corrected before being reported, because it silently mixes two
populations that should not be combined. \textbf{16 of the 20 candidates already had a baseline
proxy answer of "compliant,"} despite all 20 being dataset-labeled \texttt{ppe\_violation} --- a
known sensitivity gap in the grounding-proxy method (the same class of issue
Chapters~\ref{chap:day05} and~\ref{chap:day07} have already documented elsewhere), not something
this patch test is designed to measure. For those 16 candidates, "fooling" the model into saying
"compliant" is not a meaningful concept, because the proxy already says "compliant" before the
patch is applied at all; counting them in the denominator dilutes the real signal with samples where
the supposed "fooling" outcome was already the status quo. \textbf{Restricting to the 4 candidates
the proxy already called "violation"} --- the only 4 where a flip to "compliant" actually means
something --- \textbf{3 of 4 (75\%) flipped to "compliant" after the patch was applied.}
\texttt{09\_robustness\_test.py} now prints both numbers explicitly, the naive 15\% (3/20) and the
denominator-corrected 75\% (3/4), rather than reporting only the naive figure.

$n=4$ is far too small to generalize beyond "plausible, worth a larger run," and the four traced
cases show no clean placement-to-flip correlation, summarized in
Table~\ref{tab:day09-patchstretch}.

\begin{center}
\begin{tabular}{lp{8.6cm}c}
\toprule
\textbf{Image} & \textbf{Patch placement} & \textbf{Flipped to "compliant"?} \\
\midrule
\texttt{0000233} & off-target (torso/shoulder, not head --- a crude \texttt{\_head\_box}
approximation, no pose estimation) & yes \\
\texttt{0000641} & on-target (squarely on the head) & no \\
\texttt{0001431} & on-target & yes \\
\texttt{0001842} & on-target, scene already has many real hard-hat wearers & yes \\
\bottomrule
\end{tabular}
\captionof{table}{The 4 patch-stretch candidates for which a flip is a meaningful outcome (baseline
proxy answer already "violation"), with patch placement traced individually against the flip
outcome.}
\label{tab:day09-patchstretch}
\end{center}

Figure~\ref{fig:day09-stretch} shows the two cases that make the placement question concrete: a
well-placed patch (\texttt{0000641}) that failed to fool the model, against an off-target patch
(\texttt{0000233}) that succeeded.

\begin{figure}[h]
\centering
\begin{subcaptionblock}{0.48\linewidth}
\centering
\includegraphics[width=\linewidth]{robustness/verify/0000233_patch_before.png}
\caption{\texttt{0000233} before: baseline "violation."}
\end{subcaptionblock}
\hfill
\begin{subcaptionblock}{0.48\linewidth}
\centering
\includegraphics[width=\linewidth]{robustness/verify/0000233_patch_after.png}
\caption{\texttt{0000233} after: patch placed off-target (torso/shoulder, not head) --- flipped to
"compliant" anyway.}
\end{subcaptionblock}

\vspace{0.6em}

\begin{subcaptionblock}{0.48\linewidth}
\centering
\includegraphics[width=\linewidth]{robustness/verify/0000641_patch_before.png}
\caption{\texttt{0000641} before: baseline "violation."}
\end{subcaptionblock}
\hfill
\begin{subcaptionblock}{0.48\linewidth}
\centering
\includegraphics[width=\linewidth]{robustness/verify/0000641_patch_after.png}
\caption{\texttt{0000641} after: patch placed squarely on-target (the head) --- did \emph{not}
flip.}
\end{subcaptionblock}
\caption{The counterintuitive on-target/off-target pair: a well-placed patch (bottom) failed to
fool the model, while an off-target patch (top) succeeded.}
\label{fig:day09-stretch}
\end{figure}

A well-placed patch failed to fool the model while an off-target one succeeded --- the flip does not
cleanly track anatomically correct placement in either direction. Florence-2's "hard hat" grounding
appears generally imprecise about \emph{where} on a person the visual cue needs to be, sometimes too
lenient (an off-target patch still registers as "compliant") and sometimes not lenient enough
(an on-target patch on \texttt{0000641} did not). This is exactly the kind of result
\texttt{technical-limitations-take-and-mitigation.md} anticipated for a bounded stretch task: a real,
traceable signal, explicitly not strong enough to claim a demonstrated real-world spoofing
vulnerability from $n=4$.

\section{Standard vs. Non-Standard Choices}

\textbf{Standard}: perturbation-based robustness testing --- blur, noise, occlusion, and related
image transformations rerun through the same model and compared against a baseline --- is itself
standard practice, the same family of method behind robustness as one of the six metrics defined for
black-box XAI evaluation \citep{arreche2024exai}. Holding decoding deterministic and identical to
the logged baseline across all reruns, specifically to avoid confounding perturbation sensitivity
with Chapter~\ref{chap:day07}'s already-documented sampling-noise sensitivity, follows that
established practice rather than departing from it.

\textbf{Non-standard}: the explicit separation of answer-level and explanation-level robustness into
\textbf{two never-blended numbers} is this chapter's central methodological contribution, motivated
directly by Finding~2's traced concrete failure case (image \texttt{0000034}) rather than asserted on
principle alone. The six-metric framework's robustness definition does not by itself mandate this
split; the natural default is a single \texttt{answer\_changed} rate, and that default is exactly
what would have certified 28.3\% of stable-answer samples as fully robust despite a near-total
evidence-box relocation underneath. The denominator-correction in the patch-stretch test (restricting
to the 4 candidates for which a flip is a meaningful outcome, rather than reporting the naive 15\%
over all 20) is this chapter's second non-standard, reportable choice: catching a vacuous-denominator
issue before publishing a headline number, rather than after.

\section{Implications for the Paper}

\textbf{Findings 1 and 2, taken together, make the strongest single argument in this pilot for the
paper's central thesis that answer-level metrics are necessary but not sufficient for VLM
explainability.} Finding~1 shows that an answer-level metric can flip for a reason that has nothing
to do with the safety concept being tested (a degenerate fallback branch, not genuine sensitivity to
the occluded region). Finding~2 shows the complementary failure: an answer-level metric can stay
\emph{perfectly stable} while the explanation underneath it relocates almost completely. Read
side by side, the two findings close off both directions an answer-only evaluation could fail in ---
an answer that changes for the wrong reason, and an answer that does not change despite the
explanation collapsing --- and neither failure mode is visible from \texttt{answer\_changed} alone in
either case. \textbf{This chapter's two traced case studies, \texttt{0000341} (Finding~1) and
\texttt{0000034} (Finding~2), are the best candidates in the entire pilot for the paper's
qualitative-analysis figure}: both are visually self-explanatory before/after pairs that make an
abstract methodological point concrete in a single glance, without requiring the reader to parse an
aggregate table first.

The answer-level-versus-explanation-level split established here recurs as this pilot's central
cross-chapter theme: Chapter~\ref{chap:day10}'s Bounded Completeness analysis, Chapter~11's
cross-metric rollup, and Chapter~14's open questions for Professor Abdallah all return to this same
distinction, and this chapter is where it was first motivated by a directly-traced failure rather than
introduced as an abstract design choice. Three secondary points should also carry forward.
\textbf{First}, Finding~3's caution that rule\_2's two phrasings ground genuinely different physical
objects (a harness and a lanyard), not two wordings of one concept, should inform any future revision
of \texttt{prompts.py} before rule\_2's rewording rate is cited as evidence of model instability.
\textbf{Second}, Finding~4's confirmation that the confidence proxy is uncalibrated --- rising, not
falling, under exactly the two most disruptive perturbations --- reinforces the standing
recommendation, carried since Chapter~\ref{chap:day03}, that this pilot's "confidence" figure never
be reported as a correctness or trust signal in the paper without that caveat attached explicitly.
\textbf{Third}, the patch-stretch test's denominator-correction is worth stating as a general
methodological note in the paper's own methods section: any spoofing- or adversarial-style probe of a
binary safety classifier must first verify what fraction of its candidate set already sits on the
target side of the decision before computing a "fooling rate," or risk silently diluting a real
signal with samples for which the test's notion of "fooling" was never well-defined to begin with.
